\documentclass[10pt, letterpaper]{article}

% Packages:
\usepackage[
    ignoreheadfoot, % set margins without considering header and footer
    top=1 cm, % seperation between body and page edge from the top
    bottom=1 cm, % seperation between body and page edge from the bottom
    left=1 cm, % seperation between body and page edge from the left
    right=1 cm, % seperation between body and page edge from the right
    footskip=1 cm, % seperation between body and footer
    % showframe % for debugging
]{geometry} % for adjusting page geometry
\usepackage{titlesec} % for customizing section titles
\usepackage{tabularx} % for making tables with fixed width columns
\usepackage{array} % tabularx requires this
\usepackage[dvipsnames]{xcolor} % for coloring text
\definecolor{primaryColor}{RGB}{0, 0, 0} % define primary color
\usepackage{enumitem} % for customizing lists
\usepackage{fontawesome5} % for using icons
\usepackage{amsmath} % for math
\usepackage[
    pdftitle={John Doe's CV},
    pdfauthor={John Doe},
    pdfcreator={LaTeX with RenderCV},
    colorlinks=true,
    urlcolor=primaryColor
]{hyperref} % for links, metadata and bookmarks
\usepackage[pscoord]{eso-pic} % for floating text on the page
\usepackage{calc} % for calculating lengths
\usepackage{bookmark} % for bookmarks
\usepackage{lastpage} % for getting the total number of pages
\usepackage{changepage} % for one column entries (adjustwidth environment)
\usepackage{paracol} % for two and three column entries
\usepackage{ifthen} % for conditional statements
\usepackage{needspace} % for avoiding page brake right after the section title
\usepackage{iftex} % check if engine is pdflatex, xetex or luatex

% Ensure that generate pdf is machine readable/ATS parsable:
\ifPDFTeX
    \input{glyphtounicode}
    \pdfgentounicode=1
    \usepackage[T1]{fontenc}
    \usepackage[utf8]{inputenc}
    \usepackage{lmodern}
\fi

\usepackage{charter}

% Some settings:
\raggedright
\AtBeginEnvironment{adjustwidth}{\partopsep0pt} % remove space before adjustwidth environment
\pagestyle{empty} % no header or footer
\setcounter{secnumdepth}{0} % no section numbering
\setlength{\parindent}{0pt} % no indentation
\setlength{\topskip}{0pt} % no top skip
\setlength{\columnsep}{0.15cm} % set column seperation
\pagenumbering{gobble} % no page numbering

\titleformat{\section}{\needspace{4\baselineskip}\bfseries\large}{}{0pt}{}[\vspace{1pt}\titlerule]

\titlespacing{\section}{
    % left space:
    -1pt
}{
    % top space:
    0.3 cm
}{
    % bottom space:
    0.2 cm
} % section title spacing

\renewcommand\labelitemi{$\vcenter{\hbox{\small$\bullet$}}$} % custom bullet points
\newenvironment{highlights}{
    \begin{itemize}[
        topsep=0.10 cm,
        parsep=0.10 cm,
        partopsep=0pt,
        itemsep=0pt,
        leftmargin=0 cm + 10pt
    ]
}{
    \end{itemize}
} % new environment for highlights


\newenvironment{highlightsforbulletentries}{
    \begin{itemize}[
        topsep=0.10 cm,
        parsep=0.10 cm,
        partopsep=0pt,
        itemsep=0pt,
        leftmargin=10pt
    ]
}{
    \end{itemize}
} % new environment for highlights for bullet entries

\newenvironment{onecolentry}{
    \begin{adjustwidth}{
        0 cm + 0.00001 cm
    }{
        0 cm + 0.00001 cm
    }
}{
    \end{adjustwidth}
} % new environment for one column entries

\newenvironment{twocolentry}[2][]{
    \onecolentry
    \def\secondColumn{#2}
    \setcolumnwidth{\fill, 4.5 cm}
    \begin{paracol}{2}
}{
    \switchcolumn \raggedleft \secondColumn
    \end{paracol}
    \endonecolentry
} % new environment for two column entries

\newenvironment{threecolentry}[3][]{
    \onecolentry
    \def\thirdColumn{#3}
    \setcolumnwidth{, \fill, 4.5 cm}
    \begin{paracol}{3}
    {\raggedright #2} \switchcolumn
}{
    \switchcolumn \raggedleft \thirdColumn
    \end{paracol}
    \endonecolentry
} % new environment for three column entries

\newenvironment{header}{
    \setlength{\topsep}{0pt}\par\kern\topsep\centering\linespread{1.5}
}{
    \par\kern\topsep
} % new environment for the header

\newcommand{\placelastupdatedtext}{% \placetextbox{<horizontal pos>}{<vertical pos>}{<stuff>}
  \AddToShipoutPictureFG*{% Add <stuff> to current page foreground
    \put(
        \LenToUnit{\paperwidth-2 cm-0 cm+0.05cm},
        \LenToUnit{\paperheight-1.0 cm}
    ){\vtop{{\null}\makebox[0pt][c]{
        \small\color{gray}\textit{Last updated in September 2024}\hspace{\widthof{Last updated in September 2024}}
    }}}%
  }%
}%

% save the original href command in a new command:
\let\hrefWithoutArrow\href

% new command for external links:


\begin{document}
    \newcommand{\AND}{\unskip
        \cleaders\copy\ANDbox\hskip\wd\ANDbox
        \ignorespaces
    }
    \newsavebox\ANDbox
    \sbox\ANDbox{$|$}

    \begin{header}
        \fontsize{25 pt}{25 pt}\selectfont Avishek Bose \\
        \fontsize{10 pt}{10 pt}\selectfont Data Scientist (AI)

        \normalsize
        Location: Knoxville, TN, United States | Immigration status: Permanent resident (US)
        \\
        \mbox{\hrefWithoutArrow{mailto:avishek85.15th@gmail.com}{avishek85.15th@gmail.com}}%
        \kern 5.0 pt%
        \AND%
        \kern 5.0 pt%
        \mbox{\hrefWithoutArrow{+1(785)317-6675}{(785)317-6675}}%
        \kern 5.0 pt%
        \AND%
        \kern 5.0 pt%
        \mbox{\hrefWithoutArrow{https://linkedin.com/in/avibose}{linkedin.com/in/avibose}}%
    \end{header}

    \vspace{- 0.3 cm}
    \section{Summary}
    Applied AI researcher and engineer with 8+ years of experience designing, fine-tuning, and evaluating LLMs, RAG pipelines, and generative AI systems at enterprise scale. Experienced across the full model lifecycle --- data curation, experimentation, evaluation, and production deployment. Author of 25 peer-reviewed papers (IPDPS, IEEE BigData, ICMLA, ISI); co-designed a graduate-level AI/ML course at Kansas State University.

    \vspace{-1mm}
    \section{Work Experience}

    \begin{twocolentry}{Jun 2025 -- Present}
        \textbf{Platform Engineer}
    \end{twocolentry}
    \begin{twocolentry}{Knoxville, TN, USA}
       Tegna Inc.
    \end{twocolentry}

    \begin{onecolentry}
        \begin{highlights}
            \item Designed and deployed a production RAG platform using Azure Functions, vector databases, and self-hosted LLMs, serving 100K+ daily content discovery requests for newsroom teams.
            \item Developed hybrid search workflows combining dense semantic retrieval with keyword-based ranking, improving content retrieval precision by $\sim$35\% over keyword-only baselines.
            \item Established automated model evaluation pipelines using RAGAs and DeepEval, benchmarking 5+ production RAG components across retrieval accuracy, generation quality, and latency dimensions.
        \end{highlights}
    \end{onecolentry}

    \vspace{1mm}

    \begin{twocolentry}{Jan 2023 -- Jun 2025}
        \textbf{Postdoctoral Research Associate}
    \end{twocolentry}
    \begin{twocolentry}{Oak Ridge, TN, USA}
        Critical Infrastructure Resilience group, Geospatial Science and Human Security Division, Oak Ridge National Laboratory (ORNL)
    \end{twocolentry}
    \vspace{0.10 cm}
    \begin{onecolentry}
        \begin{highlights}
            \item Developed LLM-based NLP systems for enterprise knowledge retrieval and incident classification across 50K+ DOE/ORNL safety reports, enabling structured analysis at scale.
            \item Fine-tuned Transformer models (BERT, LLaMA) using LoRA for domain-specific classification, achieving $\sim$15\% accuracy improvement with inference latency under 100ms on Azure.
            \item Trained MatGPT on Frontier --- the world's first Exascale supercomputer --- for materials science knowledge extraction and semantic retrieval over 2M+ scientific documents.
            \item Designed scalable ML pipelines integrating embeddings, vector databases, and Azure-hosted endpoints, maintaining 99.9\%+ availability for enterprise AI services.
        \end{highlights}
    \end{onecolentry}

    \vspace{1mm}

    \begin{twocolentry}{Jun 2018 -- Dec 2022}
        \textbf{Research Assistant}
    \end{twocolentry}
    \begin{twocolentry}{Manhattan, Kansas, USA}
       KDD Lab, Department of Computer Science, Kansas State University
    \end{twocolentry}
    \vspace{0.10 cm}
    \begin{onecolentry}
        \begin{highlights}
            \item Built and evaluated NLP and graph-based ML models to extract, classify, and rank cyber-threat events from 10M+ social media posts, supporting real-time intelligence analysis.
            \item Developed GNN-based models to predict HPC job execution outcomes from cluster logs, achieving $\sim$89\% classification accuracy to support system-level resource planning.
            \item Designed an abstractive summarization pipeline using generative sequence models to process 1K+ disaster incident reports per batch, producing structured situational summaries.
        \end{highlights}
    \end{onecolentry}

    \vspace{0.1 cm}
    \begin{twocolentry}{Jan 2015 -- Dec 2016}
        \textbf{Software Engineer}
    \end{twocolentry}
    \begin{twocolentry}{Dhaka, Bangladesh}
       ICT Division, Dutch Bangla Bank
    \end{twocolentry}

    \begin{onecolentry}
        \begin{highlights}
            \item Developed backend data processing and reporting systems handling 100K+ daily banking transactions using SQL and enterprise software platforms.
        \end{highlights}
    \end{onecolentry}

    \section{Publications}
        \begin{onecolentry}
            3 journal articles and 22 peer-reviewed conference papers (IPDPS, IEEE BigData, ICMLA, ISI). \hrefWithoutArrow{https://scholar.google.com/citations?user=y5ndUuMAAAAJ\&hl=en\&oi=ao}{Google Scholar}
        \end{onecolentry}

\section{Education}

        \begin{twocolentry}{
            August 2017 -- Dec 2022
        }
            \textbf{Kansas State University}, Ph.D. in Computer Science\end{twocolentry}

        \begin{twocolentry}{
            January 2014 -- Dec 2015
        }
            \textbf{University of Dhaka}, M.S. in Computer Science and Engineering\end{twocolentry}

        \begin{twocolentry}{
            January 2009 -- Aug 2013
        }
            \textbf{University of Dhaka}, B.Sc. in Computer Science and Engineering\end{twocolentry}

\vspace{-1mm}
\section{Technical Skills}

        \paragraph{AI/ML frameworks and tools:} PyTorch, LangChain, LlamaIndex, Hugging Face, Scikit-learn, MLflow, PyTorch Geometric
        \vspace{-4mm}
        \paragraph{AI/ML:} LLM fine-tuning, Retrieval-Augmented Generation (RAG), Vector databases, Hybrid search, Agentic AI pipelines, Prompt tuning, Distillation, Quantization, Graph Neural Networks
        \vspace{-4mm}
        \paragraph{Cloud \& Infrastructure:} Azure, Docker, Terraform
        \vspace{-4mm}
        \paragraph{Programming Languages:} Python, C\#, C


\end{document}
